\scanpage{390}\mspage{304}

En effet\unsure{,|.} la vérification des propriétés a) à d), pour le système inductif canonique, de $\underset{\longrightarrow}{\lim}\simeq x$, est immédiate.

La proposition précédente ne dit pas si on peut prendre pour $J$ un ens. ordonné. Dans ce cas, la condition c) signifie que pour \unsure{tte} partie $J'$ de $J$ \unsure{grde|grd} devant $\pi_0$ et de cardinal $\leq\pi$, $J$ admet une borne supérieure $i_J$ dans $J$, et la condition d) que l'on a alors

\[
x_{i_J}\xleftarrow{\sim}\underset{i\in J'}{\underset{\longleftarrow}{\lim}}\,x_i.
\]

L'idée naturelle, pour construire un système inductif $\pi_0$-$\pi$-adapté $\inserted{(x_j)_{j\in\unsure{J|I}}}$, avec $I$ ordonné, en partant de celui qu'on a construit, avec $J=M_\pi/x$, est de prendre pour $I$ l'ens., ordonné par inclusion, des sous-catégories $\inserted{J'}$ de \deleted{$\unsure{J_0}$}$\inserted{J_0}\inserted{\text{ (sous-catégorie }\deleted{\text{\unsure{pleine}}}\text{ pleine petite de }J\text{ équivalente à }J\text{)}}$, qui sont de cardinal $\leq\pi$ et \unsure{grdes|grde} devant $\pi_0$. \inserted{(Comparer (4.7.3.3) \unsure{page|pge} 165)} Cet ens. ordonné est bel et bien \inserted{(petit et)} \uline{grd devant} $\pi$, si on suppose que $\pi^{\pi_0}=\pi$ (cf dém. 4.7.5 p.~168--\unsure{17|170})\unsure{--|.}

% archive sheet number: 389
